\chapter{„Luther oder Arndt!“}
\quellseite{1}

\begin{center}
(Eingesandt von Pastor Kalb.)
\end{center}

In dem 4. Quartalheft der Guericke- und Rudelbach’schen Zeitschrift heißt es in einer Ströbel’schen
Kritik des 12. Jahrgangs des „Lutheraners“, nachdem daselbst von pag. 724–734 der meisten Aufsätze
dieser Zeitschrift ausführlich und lobend gedacht wurde, weiter folgendermaßen:

\begin{quote}
„Hiermit glaube ich, den ‚Lutheraner‘ nach allen Lichtseiten vorgeführt zu haben. Er zeugt
kräftig 1) wider die Herren Gebrüder Papa und Apap, 2) wider die vielnamige Sippschaft der Schwarm- und
Freigeister (Wiedertäufer, Sakramentirer, Supranatural- und Rationalisten, Pan- und Atheisten, Union-
und Indifferentisten), deren Weg ‚vom Fortschritt ins Viehthum‘ führt, wider die falschen Brüder
von der Kirchen- und Amtstreiberei. Sonach steht er ganz auf dem Boden unserer evangelisch-lutherischen
Vorfahren? Doch wohl nicht ganz. Und worin liegt denn der Unterschied? Wer Geister zu unterscheiden versteht,
wird ihn wohl finden; zum Ueberfluß ist er sogar an mehreren Stellen personificirt anzutreffen. So
unter Anderm gleich in folgendem Satze des Vorworts: ‚Ohne Zagen schließen wir uns dem großen,
in viel tausend Gliedern bereits im Himmel angekommenen, Zuge an, in welchem, einen Martin Luther
an der Spitze, einst ein Johannes Brenz, ein Martin Chemnitz, ein Jakob Andreä, ein Johann und Paul
Gerhardt, ein Johann Arndt, ein Heinrich Müller und tausend und aber tausend andere treue Lehrer, zum
Theil Lehrer der ganzen Christenheit, gestanden haben, treu der Fahne des reinen Bekenntnisses unserer
Kirche folgend.‘ Aehnlich S. 13. Ganz handgreiflich auch noch S. 32.: ‚Johann Arndt’s Erklärung,
welches der Zweck seiner Bücher vom wahren Christenthum sei und wie er dieselben verstanden wissen
wolle. So schreibt der gottselige Arndt zum Schluß des zweiten Buchs vom wahren Christenthum: „Zum
Beschluß muß ich den christliebenden Leser noch etlicher Punkte freundlich erinnern: daß ich keinen
andern Finem, Ziel und Zweck habe und suche in diesen meinen Büchern, denn daß neben und mit dieser
reinen Religion und Glaubensbekenntniß, so in den Kirchen der augsburgischen
\quellseite{2}
Confession gehalten und in Formula Concordiæ wiederholet ist (zu welcher ich mich
auch mit Herzen und Munde bekenne, will auch, daß diese Schriften nicht anders, denn nach derselben
sollen verstanden werden), auch das heilige christliche Leben möge fortgepflanzt werden.“ Möchten sich
das diejenigen merken, die jetzt so oft über dem theuren Arndt inquisitorisch zu Gericht sitzen und sich
vielleicht nie so aufrichtig und herzlich zu den symbolischen Büchern unserer Kirche bekannt und nach
denselben ihre Sachen haben gerichtet wissen wollen. Möchten dies aber auch diejenigen merken, welche
jeden Eifer für reine Lehre und jedes ernste Festhalten am kirchlichen Bekenntniß für ein Zeichen
von todter Orthodoxie ansehen und erklären, während sie doch einem Arndt und anderen verstorbenen
gottseligen Theologen Gräber bauen, die mit einem Eifer und einer Treue an der reinen Lehre und Kirche
gehangen und falsche Lehre bekämpft und verdammt haben, wie sich beides jetzt kaum irgendwo findet.“

\end{quote}

Diese Stellen, namentlich die letzteren, impliciren eine tadelnde Censur gegen die ganze frühere
glaubenstreue evangelisch-lutherische Kirche, wenn sie auch zunächst bloß gegen die Buffaloer
Kryptopapisterei gerichtet zu sein scheinen.

Denn diese hatte unter anderen in dem 2ten Synodalbriefe (von 1850) S. 49 folgendes Urtheil über Arndt
abgegeben: „Bücher, wie die von Ph. J. Spener, A. H. Francke, Schubert in Potsdam, Fuhrmann in Schlesien,
Woltersdorf, Bogatzky, Rambach, Rieger, Brastberger, Richter in Barmen (Hausbibel; — auch gehört dahin die
Berleburger Bibelerklärung), und andere, lehren nur wenig oder nichts von den heiligen Sacramenten und der
heiligen Absolution und Amt der Schlüssel. So ist’s auch mit dem Hofacker’schen Predigtbuche. Diesen
Mangel trägt auch leider das beliebte Arndt’s wahres Christenthum an sich. Erdmann Neumeister’s Urtheil
über Arndt’s wahres Christenthum in seiner geistlichen Bibliothek S. 860. lautet also: Der liebe Mann hat
aus Versehen und Unwissenheit Anlaß zu vielem Streit und Mißhelligkeit gegeben, durch ein Buch, welches
er vom wahren Christenthum geschrieben. Es ist nicht zu leugnen, daß solche Redensarten darin vorkommen,
welche mit dem Fürbilde der heilsamen Lehre und Aehnlichkeit des Glaubens nicht bestehen können. Daher
sie auch von den Schwärmern gar begierig ergriffen und gemißbraucht werden. Doch ist Arndt’s Herz
und Meinung reiner gewesen, als seine Feder. Maßen er nicht aus Vorsatz, sondern aus Unvorsichtigkeit,
wozu ihn mystische und schwärmerische Schriften verleitet, anstößig geschrieben hat. Und ebendaselbst
S. 1058.: Es sind Männer, welche ein und das andere daran auszusetzen haben; wiederum andere, welche
es gänzlich verwerfen; es fehlt aber auch an solchen nicht, die es vertheidigen. Besonders streichen
es die Schwärmer und Pietisten gewaltig heraus und mangelt wenig, daß sie es nicht gar canonisirt und
der heiligen Schrift gleich gemacht haben. . . Ich unterscheide Arndt’en und Arndt’s Buch, Arndt’s
Herz und Arndt’s Feder. Er war ein frommer, unsträflicher, gottseeliger Theologus, und ist ihm wohl
nie ins Herz gekommen, etwas Irriges zu schreiben. Er hat, ehe er ein Haar breit von der
\quellseite{3}
Wahrheit abweichen wollte, sich lieber von den Calvinischen verfolgen und aussagen
lassen. Da er auch wegen verdächtiger Redensarten erinnert wurde, nahm er’s an mit sanftmüthigem
Geiste, gestund die Fehler, und protestirte mündlich und schriftlich, daß er seine Bücher anders nicht,
als nach der heiligen Schrift und unsern symbolischen Kirchenbüchern verstanden haben wollte. Allein
der liebe Mann hatte auf Universitäten nicht Zeit gehabt, einen festen Grund in der Gottseligkeit zu
legen. Denn er hatte erst Medicinam studirt, und des theologischen Pfuschers Paracelsi Schriften fleißig
gelesen, mithin die Liebe zu andern mystischen Büchern eingesogen. Weil denn diese unter schönen, hohen,
tiefsinnigen und neu gemachten Worten viel Unreinigkeit in sich stecken haben, und auch schwärmerische
Schriften mit dergleichen losem Kalf übertünchet sind, so war bei dem guten Arndtio, da er sich
hernach auf die Theologie applicirte, manches kleben geblieben. Was denn nun seine Bücher vom wahren
Christenthum insonderheit anbelangt, so leugnete er selber nicht, daß unvorsichtige Redensarten darin
eingeschlichen, und gesetzt, daß er zwölf Capitel aus dem General-Schwärmer Weigelio genommen, von dem
er aber nicht gewußt, daß er im bemeldeten Buche fast wenig vom Glauben saget, ohne welchen doch kein
wahres Christenthum sein kann, sondern es am meisten und fast durchgehends auf ein unsträfliches Leben
bauet.... Doch dieses möchte noch hingehen, wenn nur das Buch von ungesunden Redensarten rein wäre. Ob
man wohl die meisten zur Noth entschuldigen und ihnen mit einer bequemen Auslegung abhelfen kann, so
finden sich dennoch welche, die mit der Aehnlichkeit des Glaubens und dem Fürbilde der heilsamen Lehre
nicht bestehen, und daher unmöglich gut geheißen werden können. Meinen wir nun nicht, daß bedächtige
und vorsichtige Theologi Ursache gehabt, wegen dieses Buches Erinnerung zu thun? So sage ich demnach,
wenn es Leute von geübten Sinnen lesen, können sie es wohl brauchen, und das, was ihnen anstößig
fällt, leicht aus dem Weg räumen; bin aber auch gänzlich bei mir selbst beredet, weil es von einem
schweren Stylo, und mit vielen hohen und dunkeln Redensarten angefüllet ist, daß es die Einfältigen
nicht verstehen noch den vermeinten Nutzen daraus schöpfen werden.“ — Bei aller Apathie gegen
den Buffaloer Unfug kann ich doch in obiger Aeußerung nur die allgemeine Stimme der echtlutherischen
Vorzeit, und in den vorher angeführten Erklärungen des ‚Lutheraners‘ ein Abweichen davon, ein
stillschweigendes Protestiren dagegen erkennen. Es liegt hier unbestreitbar eine principielle Differenz
vor, die ich aber aus schonender Rücksicht bloß personificirend darlegen will. Nach dem ‚Lutheraner‘
muß gesagt werden: Luther und Arndt! Von unsern gläubigen Vorfahren aber wird alternativ gefragt:
Luther? oder Arndt? Von ihnen ist Arndt nie unter die normal ,treuen‘ Kirchenlehrer gezählt worden;
daß man’s heute, und nicht bloß in Amerika, sondern vielleicht noch viel mehr in Deutschland, thut,
zeugt eben von der Disharmonie, in der man sich mit der kirchlichen Vorzeit befindet. Arndt ist kein
Lehrer im Sinn und Geist der Reformation. Nun behauptet er freilich zur Ablehnung
\quellseite{4}
dieses Vorwurfs in der Vorrede über das erste Buch vom wahren Christenthum:
‚Es ist der erste Artikel von der Rechtfertigung des Glaubens in diesem, sonderlich aber im andern
Buche also geschärfet, und so hoch getrieben, als es immer möglich!‘ Wohl! Aber statt daß nun
diese gewaltige evangelische Hauptlehre den ganzen übrigen Inhalt des ‚wahren Christenthums‘
durchdringen und regieren sollte, steht sie einflußlos, isolirt da, und Gesetz und Mystik führen das
Regiment; weshalb Neumeister mit Recht bemerkt, Arndt sage ‚fast menia vom Glauben‘. Denn sowie
er den Artikel von der Rechtfertigung episodisch abgehandelt hat, eilt er, fröhlich wie Einer, der
ein unliebsames Werk ein für allemal abgemacht, wieder an sein nur momentan unterbrochenes Geschäft,
Mosen und Weigeln zu treiben. Aber Arndt will doch, nach dem obigen Citate des ‚Lutheraners‘, seine
ganze Schriftstellerei nach den symbolischen Büchern verstanden wissen? Es ist jenes Citat nicht die
einzige, nicht einmal die bestimmteste Stelle, wo er das verlangt. Er spricht sich gleichermaßen noch
in der Vorrede über das 4te Buch aus und in der Vorrede zum ersten Buch heißt es: ‚Ich protestire
hiermit, daß ich dies Büchlein, gleichwie in allen andern Artikeln und Puncten, also auch in articulo
de libero arbitrio, justificatione peccatoris coram Deo nicht anders, denn nach dem Verstande librorum
symbolicorum Augustanae confessionis, als da sind die erste unveränderte Augsburgerische Confession,
Apologia, Schmalkaldische Artikel, beide Catechismi Lutheri und Formula Concordiae, will verstanden
haben.‘ Ich bekenne offen, daß ich eine solche Art der Berufung auf die symbolischen Bücher, die mir
bei keinem unserer glaubensstreuen Theologen je vorgekommen ist, gar nicht verstehe, dieselbe vielmehr
für das lächerlichste und doch gefährlichste Ding halte, — für das lächerlichste, weil es sich,
genau betrachtet, in dem Köhlerglaubensringe herumtreibt (ich glaube, was das Kirchenbekenntniß
glaubt, und dieses glaubt, was ich glaube, — was aber das Bekenntniß glaubt, kann man nicht aus
meinen Worten, und was ich glaube, nicht aus den Worten des Bekenntnisses ersehen; — denn so steht
doch thatsächlich das Verhältnis zwischen den symbolischen Büchern und den Büchern Arndt’s ‚vom
wahren Christenthum‘), — für das gefährlichste, weil es allen Irrthümern den Eingang in die Kirche
öffnet; denn mit gleichem Rechte wie Arndt könnten auch die Rationalisten, Unionisten, Pantheisten
verlangen, daß man ihre Schriften der Augsburger Confession u. s. w. gemäß verstehe; würde damit
nicht die Irrlehre, ohne nur einen Buchstaben zu widerrufen, mit einem Schlage für rein und orthodox
erklärt? Ich wenigstens kann mich nicht überzeugen, daß die doctrinelle Reinheit der symbolischen
Bücher die mangelnde Orthodoxie der Kirchenlehrer vertreten und ersetzen könne. Im Gegentheil glaube
ich, solchen Vorstellungen sei mit allem Nachdruck entgegen zu treten, namentlich in Deutschland, wo
eine evangelisch und lutherisch sein wollende, aber in der That nicht über Mosen und Weigeln (oder wie
man die aus eigenen Geistestiefen oder Seichtigkeiten schöpfende, fromm-speculative Beschauung sonst
\quellseite{5}
nennen mag) hinauskommende Welt- und Lebensbetrachtung, statt bekenntnismäßig zu
lehren, es vorzieht, ihre Bekenntniswidrigkeiten symbolisch deuten zu lassen. Davon weiter zu sprechen
ist hier nicht der Ort; unserm amerikanischen ‚Lutheraner‘ gereicht es aber auch nicht zum Vortheil,
daß er Luther und Arndt Arm in Arm, als verstünde sich das von selbst, mit einander wandeln läßt. Er
hat dadurch neben der ursprünglich reformatorischen auch noch eine andere Physiognomie erhalten: jene
erbauliche Vielgeschäftigkeit, die, wenn je, so gerade jetzt Bedenken erregen muß; — sehen wir doch,
daß Leute, wie Bunsen, die am christlichen Glauben Schiffbruch gelitten haben, nicht Schmähworte genug
gegen die evangelisch-lutherische Kirche ob ihres Mangels an gewerblicher Gottseligkeit (1 Timoth. 6,
5.) zu finden wissen. So betrübt es uns denn, daß auch der ‚Lutheraner‘ des frommen Zeitgeistes
Gaukelsack, die Mission, handhabt, ohne zu bedenken, daß die beiden stammverwandten Namen Missa
und Missio die fruchtbaren Mütter aller kirchlichen \unsicher{Werkrümpel} sind. Die Messe hat’s Alles regiert,
klagte Luther, — die Mission regiert’s Alles, klagt die tägliche Erfahrung. — Eine oberflächliche
Betrachtungsweise könnte freilich schnell sagen, vom Standpunkte des ‚Lutheraners‘ sei es sehr leicht,
über die gerügten Punkte die zufriedenstellendste Erklärung zu geben. Ja, wenn nur erklärende Worte den
Lauf der Dinge aufzuhalten oder zu ändern vermöchten! Sehe nur jeder selbst, wie sich das Evangelium seit
hundert Jahren Bahn gemacht hat. Etwa durch gütliche Versicherungen und Transactionen? Nein, unter harten
Kämpfen zuerst mit der Verdummung der Aufklärer, dann mit der pietistischen Frommthuerei, hernach mit
dem Unionsgreuel, und endlich mit der Kirchen- und Amtstreiterei. Und überall ließen es die Gegner an
glänzenden Worten nicht fehlen, um sich als ‚Brüder‘ der Evangelisch-Lutherischen zu dokumentiren,
und klagten laut über deren ungefugen Sinn, wenn der Gang der Ereignisse die Trennung des nicht zusammen
Gehörigen consolidirte. So und nicht anders wird es auch hier in der erwähnten Angelegenheit kommen. Der
noch nicht ausgefochtene Amts- und Kirchenstreit verdeckt vor der Hand noch die wahre Lage der Dinge, aber
hinter ihm lauscht bereits erkenntlich genug die Frage: Luther? oder Arndt? (man formuli[e]re sie, wie man
wolle) hervor. Ihre Beantwortung wird keine friedliche sein; die Zahl derer, welche den unerschrockenen
Mitstreiter gegen Heidenthum, Pfaffenthum und Cäsarenthum neben den Reformator stellen, ist in beiden
Hemisphären nicht gering und die dahin gehörenden Persönlichkeiten sind achtunggebietend; schmerzlicher
und schädlicher als in den bisherigen Kämpfen wird der Riß werden, den jene Frage unter den Bekennern
der Augsb. Conf. verursachen wird und den aufzuhalten doch in keines Menschen Macht steht. Auch wäre
es für die Zukunft der evangelisch-lutherischen Kirche nichts weniger als heilsam, wenn jene Frage ohne
gründlichen Austrag bliebe; so gewiß es unserer kirchlichen Zukunft nur zum Segen gereichen wird, daß
sie geläutert durch die rationalistischen, pietistischen, unionistischen und Amts-Kämpfe aus der Gegenwart
\quellseite{6}
aufsteigt. Hüten wir uns aber, das allerdings schon unter der Asche glimmende Feuer vor
der Zeit anzufachen; die leichte Decke, die es gegenwärtig noch für viele Augen verhüllt, wird ohnehin
rasch genug von den in unserm Jahrhundert so scharf wehenden kirchlichen Winden hinweggeblasen sein. Auf
diesen Ausgang schon in Zeiten aufmerksam zu machen, nicht aber ihn herbeiführen zu helfen, dürfte
mehr als manches Andere die Pflicht der aufrichtigen, wohlmeinenden evangelischen Theologie sein.“ —

So weit Ströbel. Das ist sein Judicium — nach seiner Art. Wer ihn kennt, wird sagen: Es ist der
Alte. Jedes Wort ein Pfund. Andere möchten vielleicht in etwas anderer Weise ihren Herzen Luft machen,
etwa also: „Was fragt ihr viel nach dem Urtheil eines Solchen, der mit Worten drein wirft, als wären
und müßten es lauter reformatorische Felsflöße sein?“ „Laßt den Eisernen doch eisern!“ „Laßt
ihn den Propheten spielen; es wird doch bald ausgespielt sein und wer zuletzt lacht, lacht am besten.“
„Wenn das Zeichen nicht kommt, das er weissagt, wenn zur Zeit des Risses es auch von ihm heißt: ‚Ich
habe mir lassen überbleiben‘ — auch über’m Ocean — sieben tausend Mann, die nicht haben ihre
Kniee gebeugt!‘ dann werden ihm die Augen wohl aufgehen!“ Doch offen und ehrlich — wir gehören
nicht zu Solchen. Uns lacht das Herz im Leibe, so oft wir’s lutherisch rauschen und sausen hören,
wenn es auch zu Zeiten scharf bläßt und Windstöße gibt. Wir wollen jenen, wir können ihnen nicht
folgen. „Den Geist dämpfet nicht.“ Jeder handle und zeuge, „nachdem Gott ausgetheilt hat das Maß
des Glaubens.“ Wer zum Eisernen berufen ist, der eisere — ja, eisere heilig und recht.

Wir wollen uns aber doch rechtfertigen; freilich nur kurz. Viele Worte sind hier nicht Noth. Hat doch
unser verehrter Kritiker deren genug vor uns gesprochen. Und noch mehr: wir sind nicht von denen, die
da meinen, „erklärende Worte sollten den Lauf der Dinge aufhalten oder — ändern.“ Ist das arme
Amerika, das nach einer deutschen, gelehrten Vision vom dominirenden calvinischen Geist gänzlich und
hoffnungs- und rettungslos verschlungen ist, auch sonst übel genug bestellt, so mangelt ihm doch noch
unseres Wissens der zweideutige Ruhm, im Produciren großer Wort- und Maulhelden besonders fruchtbar zu
sein. Darum geschehe die Rechtfertigung nur — kurz; statt der Worte — Thaten.

„Wer fordert solches von deinen Händen?“ — Zu dieser kurzen Rechtfertigung und thatkräftigen
„Erklärung“, die gar nicht intendirt, ins Rad der flüchtigen Zeit operirend eingreifen zu wollen,
sondern nur in tiefster Bescheidenheit, obwohl mit lutherischer Ehrlichkeit und Gradheit Thaten wider
Worte will auftreten lassen, hat der Unterzeichnete desto mehr Lust, weil es sich hiebei nicht um eine
Schlacht auf dem Feld der Wissenschaft handelt, sondern nur um Darstellung gewisser Facta, und zweitens,
weil der Geringste aus dem entgegengesetzten Feldlager in seiner heiligen Waffenrüstung wider einen
Helden lutherischer Kernhaftigkeit, wider einen Elias-Eiferer im Goliaths-Costüme
\quellseite{7}
durch seine 5 Schleudersteine doch vielleicht ein größeres Gewicht in
die Waagschaale legt, als wenn einer der Besten hervor und auf den Plan treten müßte. In Kurzem, wer
es nicht weiß, der wisse es nun: „ein schlechter, einfältiger Prediger“ des Evangeliums ist es
(mit Luther zu reden), der da schreibt, „erklärt“, opponirt. Das sei genug zu meiner persönlichen
Rechtfertigung; nun zu der der Sache.

Ströbel’s Angriff ist scharf, und soll ich’s genau sagen: beißend „mit schonender Rücksicht“, oder
soll ich’s auf’s gelindeste ausdrücken: etwas ungerecht in Folge der scharfsichtigen Kurzsichtigkeit. Wer
kann es aber auch verlangen, daß das lutherisch-scharfsichtige Auge eines Elias-Eiferers bis über
den Ocean Alles am Kirchenhimmel, die Sterne, die Wolken und Wölkchen bis aufs Kleinste haarscharf
distinguiere? Darum nur — Geduld, wenn auch die Waffen ernstlich klirren. Doch wohl verstanden! wir
gedenken nicht peccavi zu sagen, wenn wir versuchen, einem solchen Gegner entgegen zu geben. Mit einem
Solchen ist nicht viel zu scherzen. Darum gelte es! Frisch und frank zur ersten Schlacht wider den
Str.’schen Angriff.

Wir sagen: Die specifisch Arndt’sche Richtung ist so wenig die unsrige, daß wir vielmehr Str. im
sachlichen Urtheil über dieselbe von Herzen beistimmen. Kurz und gut: Das sind auch unsere Sätze,
unsere Bekenntnisse — mit Str. —: „Arndt ist kein Lehrer im Sinn und Geist der Reformation.“ Der
Artikel von der Rechtfertigung des Glaubens, „diese gewaltige, evangelische Hauptlehre, durchdringt und
regiert den ganzen übrigen Inhalt des ‚wahren Christenthums‘ nicht so, wie es sein sollte.“ Ja, um
auf’s Strengste und Schärfste mit Str. zu reden: „Gesetz und Mystik führen da das Regiment.“ Das
sind auch unsere Sätze — mit Neumeister —: „unvorsichtige und ungesunde Redensarten haben sich“
in dies Buch „eingeschlichen, die mit der Aehnlichkeit des Glaubens und dem Fürbild der heilsamen
Lehre nicht bestehen“; „er baut das wahre Christenthum in dem Buch meistens und fast durchgehends
auf ein unsträfliches Leben.“ — Und mit Pastor Grabau: „Den Mangel trägt leider auch Arndt’s
‚wahres Christenthum‘, daß es „wenig oder nichts von den heiligen Sacramenten und der heiligen
Absolution und dem Amt der Schlüssel lehrt.“

Weiter bekennen wir: 1.) daß die specifisch Arndt’sche Richtung, sonderlich nach den 6 Büchern vom
wahren Christenthum, die hier um mehrerer Gründe willen vor Allem in Betracht kommen, eine ungesunde ist,
das heißt: So einfältige Christenleute, um in gottseeliger Erkenntniß zu wachsen, im Glauben stark
und im Leben frommer zu werden, ausschließlich Arndt’s 6 Bücher vom wahren Christenthum studieren
und zu ihrer Erbauung lesen würden, so würden sie dadurch keineswegs \unsicher{ferngesunde}, fröhlichgläubige
und echt kirchlich-lutherischgesinnte Bekenner der Wahrheit werden. 2.) Die Arndt’sche Richtung ist
nach unserm geringen Dafürhalten in der That die Wurzel des Pietismus. In ihr sind darum dem Keime nach
alle die verderblichen Früchte derselben verborgen und verschlossen. 3.) Praktisch
\quellseite{8}
betrachtet, scheint uns das das Gefährlichste zu sein, daß in solchem Buch viele
Schriftstellen, die der theure Dr. Luther mit Recht also handelt, dadurch „des Glaubens Trost“
zu stärken, aufs Leben und die Lebensgerechtigkeit gezogen werden. So wird eine ungesunde Exegese
für das Volk befördert; ja, so wird für dasselbe die specifisch Arndt’sche Richtung in die liebe
Bibel selbst hineingetragen. 4.) Wider die, so da sagen: Arndt’s Weise finde ihre Berechtigung in dem
seiner Zeit so nothwendigen Kampf wider die todte Orthodoxie, behaupten wir, daß Dr. Luther die Lehre
der Lebensgerechtigkeit viel schärfer und gewaltiger treibt, als es Arndt immer möglich gewesen, a)
weil solche Lehre und Vermahnung lebendig und recht schriftgemäß aus dem Grund des Glaubens fließet,
und b) er die Nothwendigkeit, das Gesetz Gottes nach seinem dritten Brauch unaufhörlich zu handeln,
eben so gut, ja vielmehr ohne Zweifel besser, als Arndt, erkannte. Hierzu erinnere man sich nur des
seligen und hohen Streites des theuren Gottesmannes Dr. Luther wider die Antinomer und frage sich,
wer die Lehre „vom Beruf“ in Predigten und sonstigen Schriften fleißiger und praktischer treibt,
als eben der alte auch in dem Stück unübertreffliche Reformator. 5.) Sagen wir auch von ganzem Herzen
mit Neumeister von dem Verfasser des „wahren Christenthums“: „er war ein frommer, unsträflicher,
gottseliger Theologus, und ist ihm wohl nie ins Herz gekommen, etwas Irriges zu schreiben“; rühmen
wir auch gebührend, daß „er, ehe er ein Haarbreit von der Wahrheit abweichen wollte, sich lieber
von den Calvinischen hat verfolgen und ausjagen lassen“; ja stellen wir ihn auch hinsichtlich seiner
bekenntnißtreuen Demuth den „Buffaloer Unfug“-Treibern, „den Kryptopapisten“ als Exempel
und Muster dar (wie es im „Lutheraner“ geschehen), da Grabau aus Neumeisters Mund also schreibt:
„Da er auch wegen verdächtiger Redensarten erinnert wurde, nahm er’s an mit sanftmüthigem Geiste,
gestund die Fehler (!!!)“: so kann doch das auch nicht verschwiegen werden — jeder, der mit Fleiß
anfängt, Luther zu studiren, muß es ja mit uns bekennen —, daß die ganze Lehrweise Arndt’s
nicht mit der guten, alten, kirchlich-lutherischen übereinstimmt. Seine Phrases und Modi loquendi
sind andere. Ihm fehlt die Kirchensprache. So ist es nach unserer unmaßgeblichen Meinung auch nicht
ungerecht, die große Verbreitung des Arndt’schen „wahren Christenthums“ wenigstens zum Theil von
der gesetzlichen, dem natürlichen Menschen freilich besser behagenden Richtung herzuleiten, obwohl man
das Haschen der Secten und Schwärmer nach Arndt’s 6 Büchern füglich also mildernd deuten kann, daß
solche Hummeln nur das Ihre im Arndt suchen und nicht den ganzen Arndt. Sollte solches Urtheil wirklich
einem oder dem andern Leser zu schroff dünken, so fange er nur mit Ernst an, der Sache nachzuforschen;
er werde ein Schüler Dr. Luthers und setze sich zu seinen Füßen, so wird er es selbst bekennen, „es
will sich nicht recht reimen zu Arndt’s Art und Weise“; oder vielmehr umgenannt: „diese nicht zu
jenem.“ Das ist unsere erste Schlacht wider Str., indem wir
\quellseite{9}
bekennen: Sachlich hat Str., N. und Gr. Recht — NB. mit uns. Wir stimmen ihm bei. Aber
merke, nicht bloß das ist’s! — Die zweite Schlacht?

Wir thun dies nicht bloß jetzt; das war unsere Richtung längst, ja immer, seit dem Bestand unserer
Synode von Missouri, Ohio u.\,a. St., seit 12 Jahren; wir zeugten in solcher Weise schon vor Jahren in dem
von Str. später angegriffenen „Lutheraner“. Redet daher Str. auch etwas hart: er sähe hier „eine
principielle Differenz“, oder „an mehreren Stellen“ im „Lutheraner“ „sei der Unterschied“,
der da zwischen uns und „unsern lutherischen Vorfahren“ Statt fände, „personificirt anzutreffen“,
und er, nämlich Str., „wolle die principielle Differenz aus schonender Rücksicht bloß personificirend
darlegen“ — warum, mein theurer Eiferer, nach der Liebe, die da aufrichtig traf, nicht detaillirend? oder
wäre wirklich das Erstere weniger gravirend?\footnote{Hier seien uns nur folgende kurze Bemerkungen erlaubt.
1) Wir begehren, wenn es sich nicht um Leben, sondern um Lehre handelt, so wenig „Schonung“ für uns,
so wenig wir hierin Schonung haben für Andere und diesen gewähren. 2) Nichts kann schonungsloser und
darum widersprechender sein, als wenn der Kritiker selbst sagt, er wolle aus „Schonung“ das Schlimmste
nicht sagen! Gegen einen solchen verdeckten unlutherischen Angriff kann man keinen ehrlichen Kampf
eröffnen, und das ist auch die Ursache gewesen, warum wir uns bisher nicht haben überwinden und
entschließen können, auf Hrn. etc. Ströbels Kritik auch nur ein Wort zu erwidern. 3) Wüßte derselbe, wie
gerade darin die Führung besteht, damit wir hier unverdient begnadet und zu unseren hiesigen Kämpfen
gewappnet worden sind, daß wir nach jahrelangen unaussprechlich schweren Erfahrungen des Jammers und
der Noth eines gesetzlichen Pietismus das Kleinod eines wahrhaft evangelischen Lutherthums erlangt haben:
so würde er aus einzelnen mißdeutbaren Aeußerungen κατ’ ἄνθρωπον sich und Andern das Bild nicht
entworfen haben, das er von uns hat und in der besprochenen Kritik zeichnet. 4) Wenn endlich Hr. Lic.
Ströbel in Arndt den Abfall vom wahren Lutherthum, hingegen in einem Neumeister, Schlegel, Corvinus u.\,A.
die „echt lutherische Vorzeit“ und nicht auf beiden Seiten das Greisen und Weichen sieht, und wenn er im
Ernste die Missio dem gotteslästerlichen Greuel der Missa an die Seite stellt: so findet zwischen uns und
ihm allerdings eine bisher ungeahnte principielle Differenz statt, die nur von seiner Seite ausgeglichen
werden kann. Die Red.} — muß uns solches auch „etwas ungerecht“ erscheinen,
da unser verehrter Kritiker den Grund der im „Lutheraner“ angeführten Stellen wenigstens zum
Theil erkennt, daß sie nämlich „zunächst bloß gegen die Buffaloer Kryptopapisterei gerichtet
sein sollen“, so können wir doch solches wohl ertragen. \unsicher{P./N. 141, 5.} Unser „Lutheraner“ ist kein
wissenschaftlich-theologisches Blatt, sondern hat durchweg eine praktische Tendenz. Da galt es nach des
theuren Redakteurs Opinion wohl „nicht mit der Thür ins Haus zu fallen“ oder, in Ströbel’scher
Weise zu reden, „uns zu hüten, das allerdings (auch hier) schon unter der Asche glimmende Feuer vor der
Zeit anzufachen.“ Darin besteht das ganze Mysterium contradictionis nach unserm Dafürhalten. Würde
uns aber entgegnet: „Nicht also! so leichten Kaufs entgeht ihr mir nicht! In Principienfragen handelt
man anders!“ so antworten wir: Der „Lutheraner“ ist das öffentliche Organ unserer Synode. Wir
kennen ihn auch nach dem, wie er die Lehre und Praxis vertritt, wohl
\quellseite{10}
seit 14 Jahren. Hätte er es auch hier in kirchlicher Diplomatie etwas verfehlt, so gibt
ihm doch ein solches Fleckchen noch kein anderes Gesicht, nicht die von Str. ersehene, unreformatorische
„Physiognomie“, denn, das merkte nun, geneigter Leser, 5 Jahre vorher schon, ehe Str. diese geschwärzte
„Physiognomie“ ersieht, zeugt der „Lutheraner“ entschieden wider die Arndt’sche Richtung. Es
heißt im Jahrgang 7. No. 3. pag. 17 und 18 also:

„Wenn man nun zugeben muß, daß besonders bald nach der Reformationszeit die ganze Weise im Predigen
und in der Uebung des Katechismus eine andere geworden sei, so ist doch der Grund davon nicht sowohl in der
Unvollständigkeit der alten Kirchen- und Schulordnungen und noch weniger in einer gewissen Einseitigkeit
des reformatorischen Verfahrens zu suchen, denn es waren ja namentlich treffliche Muster genug vorhanden,
die Schuld lag an denen, die diesem Vorbild nicht folgten. Es kam eine andere Lehrweise auf, die mehr
der schwerfälligen Rüstung Saul’s glich, als der leichten aber sicheren Schleuder Davids; die ganze
Anlage in Predigten wurde gekünstelt und in gewisse Schulregeln eingezwängt; des geschichtlichen
Grundes wurde vergessen und an dessen Stelle trat eine trockene Ueberlieferung der Lehrartikel; über
dem nöthigen Streit gegen allerlei falsche Lehre wurde die Erbauung in der rechten Bahn versäumt; man
suchte dem überhandnehmenden Ueberdruß an dem einfachen Gottesworte durch allerlei Geschichtchen und
andere Zuthaten vorzubeugen, auch die Sprache wurde holprichter und immer unähnlicher der kraftvollen
eigenthümlichen Kirchensprache Luthers.

„Der Einfluß, welchen J. Arndt’s Richtung auf die Behandlung der Predigt und des Katechismus hatte und
die schon vom Jahre 1605 an zu datiren ist, in welchem das 1. Buch seines wahren Christenthums erschien
— dieser Einfluß ist in jenem werthvollen Aufsatze (Dr. Kliefoht’s über Predigt und Katechese in
der Vergangenheit und Gegenwart der lutherischen Kirche) gar nicht berührt worden, da doch aus dieser
Richtung die Spener’sche hervorging; die letztere aber ist so treffend und freimüthig dargestellt
worden, daß man wünschen muß, recht bald eine ähnliche Schilderung der Arndt’schen und namentlich
auch unserer neuesten Zeiten zu bekommen. Wie sich doch die Zeiten hinsichtlich des Urtheils über die
Zeit ändern! Wer vor 10 Jahren sich so über die Spener’sche Richtung geäußert hätte, der würde
ohne Zweifel von pietistisch gesinnten Lutheranern für einen in blindem Eifer befangenen Antipietisten
und darum auch für einen Antilutheraner erklärt worden sein. Gott gebe, daß auch in den beiden andern
Beziehungen bald die Zeit kommen möge, wo man endlich einmal frei herausredet.“ Ein ähnliches Zeugniß
wider die Spener’sche, und natürlich inclusive auch wider die Arndt’sche Richtung, findet sich auch
schon im 5. Jahrgang des „Lutheraners“, pag. 189, obwohl daselbst der Name Arndt’s nicht speciell
genannt wurde. — In Summa: unsere zweite Schlacht wider Str. ist ebenso siegreich als die erste. Unser
Ruhm ist der, daß wir stets den
\quellseite{11}
heilsamen Rückschritt predigten, nicht etwa zu Arndt, sondern zu Luther, daß
wir nach dem Maß des Glaubens, das uns Gott verliehen, uns stets befleißigten, allen ungesunden,
lutherisch-seinwollenden Richtungen zu widerstreben, und uns nicht scheuten, vor aller Welt als treue
Kämpfer wider den Pietismus dazustehen (nach seiner Wurzel, seinem Wesen und seinem Umfang). Darum ist,
darum war, darum bleibt mit unseres Gottes Hülfe — trotz alles Personificirens — unsere Parole nicht
etwa: Luther und Arndt! (wenn es hierin auch in der Praxis fehlen sollte), sondern einmüthigen Geistes:
„Nur Luther!“ oder schärfer ausgedrückt: „Nicht Arndt, sondern Luther!“ — Es kann ja nichts
schaden, so wir unsere Bollwerke in dieser Hinsicht selbst immer besser erkennen, um auch darin „stark
zu werden in dem Herrn und in der Macht seiner Stärke“, und treulich „ob dem Glauben zu kämpfen, der
uns vorgegeben ist.“ Das erste ist die selige Gewißheit — mit Dr. Nic. Selnecker zu reden — „daß,
so lange“ wir „in diesem Lande, in unsern Kirchen und Schulen, über diesem Bekenntniß und Erklärung,
so in dem christlichen Concordien-Buch verfaßt, halten, so lange wird auch Richtigkeit in Gottes Wort
oder in der Lehre, ohne Schwärmerei, neben anderem Segen Gottes bei uns sein und bleiben.“ Das ist es:
Von nahezu 200 Kanzeln erschallt bei uns in unserm Kreis das liebe Evangelium sonntäglich, das reine, das
volle Evangelium ohne pietistisch verklausulirte Darreichung des himmlischen Trostes, ohne die leidigen
dem Evangelischen: „Es ist Alles bereit!“ oder dem Apostolischen: „Alles ist Euer!“ nachhinkenden
Wenn’s und Aber’s. — Nein, die Applicatio bleibt voll, kräftig, ganz und unverkümmert nach dem
lutherischen General-Applicationssatz: „Wer da glaubt, der hat’s.“ „Durchdringt und regiert“
solche Lehre des Wortes, „durchdringt und regiert“ die gesunde Theilung des Wortes in Gesetz und
Evangelium, nicht nach Arndt und Spener’s Anleitung, sondern nach der, die der theuere Reformator
uns in seinen Postillen und sonstigen Schriften gibt — das sind aber die rechten kaum erschöpflichen
Fundgruben und lutherischen Bergwerkstollen, dazu jeder Prediger bei uns nicht bloß privatim, sondern
von Amtswegen, und dazu durch mehrfach wiederholte öffentliche Synodalbeschlüsse gewiesen wird —,
durchdringt und regiert Solches eine jede unserer Predigten, was gilt es?! — unsere Parole erweist
sich siegreich. So ist es auch hinsichtlich unseres Katechismus-Unterrichts. Dr. Luthers Katechismus
ist unser Handbuch. (Hier wäre auch des ausführlichen Katechismus Pastor Keyls zu gedenken als einer
echt-antipietistischen Frucht aus der Mitte unserer Synode.) Dr. Luthers Lieder füllen unsere Kirchen;
nicht die süßen, vom subjectiven Gefühlswesen überfließenden Mach- und Backwerke späterer Zeiten,
die der guten, nahrhaften, und kräftiglutherischen Hausmannskost doch so weit nachstehen, als der im
Gefühl verschwimmende, schwammartige Individualitätsglaube unserer Zeit dem kernigen, tapferen und
mannhaften Kirchen- und Bekenntnißglauben der hehren
\quellseite{12}
Reformationszeit. So fehlt es nicht und kann nicht fehlen nach solchen Prämissen: Dr. Luthers Schriften sind
nicht bloß unsere Lieblingsstudia, unsere „lieblichen Gemüßgärtlein“ und diesseitigen himmlischen Paradisi,
sie sind uns auch eine wider alle unsere Feinde wohl versehene Rüst- und Waffenkammer, aus der wir Eines nach dem
Andern hervorholen. Müssen auch unsere öffentlichen Organe, „der Lutheraner“ und „Lehre und Wehre“, oft
und vielmals von Zeitungshelden unserer Tage Spott und kleinliche Kläfferei erfahren wegen der vielen Citate aus
Dr. Luthers Schriften, so ficht uns das doch so wenig an, daß wir nicht allein unsere Rüstkammer immer fleißiger
besuchen, sondern auch in dieser Beziehung den Ausspruch eines A. Harleß in seiner Vorrede zu seinen Thesen
über „Kirche und Amt“ als einen herrlichen und köstlichen bezeichnen, wenn er dort pag. VIII. also schreibt:
„Wenn ich dann ausführlicher Luther citire, so bitte ich, das nicht als ein bloßes Citat oder Aussage eines
Andern anzusehen. Ich lasse Luther für mich reden, denn er redet besser, als ich zu reden vermöchte. Auch geht
es mit Männern solchen Berufes, daß sie das, was sie für bestimmte Zeiten gesagt haben, durch providentielle
Fügung wie für alle Geschlechter geredet zu haben scheinen. In diesem Sinne führe ich ihn ein.“ So heißt es
auch im Buch „die Stimme unserer Kirche“ gleich in der Vorrede: „Wir, die Glieder der deutsch ev.-luth. Synode
von Missouri, Ohio u.\,a. St.... sind darauf bedacht gewesen, aus Luthers Schriften besonders reichlich Auszüge zu
geben; das wird gewiß keinem Leser auffallend sein, der mit uns durch Gottes Gnade erkannt hat, daß Luther und
nur Luther“ — hier hörst du unsere Parole officiell ausgesprochen — „und nur Luther der von Gott berufene
Reformator der Kirche gewesen ist, nach dessen Namen wir uns gern nennen lassen, nicht weil wir an ihn glaubten,
sondern weil wir erkannt haben, daß die Lehre, welche er gepredigt hat, nicht seine Lehre, sondern das lautere Wort
des ewigen Gottes ist.“ In Summa: So lange bei der öffentlichen Ausrichtung unseres Lehr- und Zunamens „Gottes
Wort und Luthers Lehr“ bei uns im Schwange geht, schallt und klingt, und insonderheit die Lehre vom Glauben
nach unsers theuren Vaters und Meisters Art Alles, Alles durchfließet, durchlebet, durchgeistet und regieret; so
lange Dr. Luther bei uns lebt in seinen Schriften, in unsern Predigten, Katechesen, in unseren Liedern, bei unsern
Strafen, Ermahnungen und Trösten, in unsern heiligen Kämpfen und Streiten um des Herrn Sache: so lange ist noch
keine Gefahr, daß wir erkannter Weise unsere Parole sollten schändlich und schmählich verlassen. „Nur Luther!“

Zum Andern ist es noch nicht erwiesen, daß bei uns die Lehre von der Heiligung und den guten Werken in einem
ungesunden und unlutherischen Verhältnis zur Lehre von der Rechtfertigung und dem Glauben stehe. Sieht man uns
auch so ohnehin „einer erbaulichen Vielgeschäftigkeit“, so heißt das doch nach unserm groben und ungelehrten
Verstande noch lange nicht so viel, als sei damit Alles bewiesen; wir müssen vielmehr ernstlich bitten, das
schonte „personificirende Darlegen“ ein wenig bei Seite zu legen und lieber
\quellseite{13}
nach unsers Meisters Art etwas derbe und „unschonender“, d.\,h. lutherisch-grob zu verfahren. Denn das wird
man uns doch wohl erlauben, daß wir den kühnen, großartigen Ausfall auf die Mission, „als den Gaukelspaß
des frommen Zeitgeistes“ (was wir, beiläufig gesagt, in gut lutherischer Weise als „das Kind mit dem Bade
ausschütten“ uns zu bezeichnen verstatteten) nicht als eine hinreichende Basis zu einer solchen Beschuldigung
anzuerkennen vermögen. Darum bleiben wir noch getrost in Lehre und Praxis bei unsers Luther’s Generalproclamation
über den Zusammenhang von Glauben und Werken, der es klärlich so ausspricht:

Der aus eigenen Kräften gemachte Glaube „thut auch nichts und folgt keine Besserung hernach. Aber Glaube ist
ein göttlich Werk in uns, das uns wandelt und neu gebieret aus Gott, Joh. 1, 13., und tödtet den alten Adam,
machet uns ganz andere Menschen, von Herzen, Muth, Sinn und Kräften und bringt den Heiligen Geist mit sich. Das
ist ein lebendig, geschäftig, thätig, mächtig Ding um den Glauben, daß unmöglich ist, daß er nicht ohne
Unterlaß sollte Gutes wirken. Er fraget auch nicht, ob gute Werke zu thun sind; sondern ehe man fraget, hat er
sie gethan, und ist immer im Thun. Wer aber nicht solche Werke thut, der ist ein glaubsloser Mensch, tappet und
stehet um sich nach dem Glauben und guten Werken, und weiß weder was Glauben und gute Werke sind, wäschet und
schwatzet doch viel Worte vom Glauben und guten Werken. Glaube ist eine lebendige, erwogene Zuversicht auf Gottes
Gnade, so gewiß, daß er tausendmal darüber erschröbe. Und solche Zuversicht und Erkenntniß göttlicher Gnade
macht fröhlich, trotzig und lustig gegen Gott und alle Creaturen: welches der Heilige Geist thut im Glauben. Daher
der Mensch ohne Zwang willig und lustig wird, Jedermann Gutes zu thun, Jedermann zu dienen, allerlei zu leiden,
Gott zu Liebe und zu Lob, der ihm solche Gnade erzeiget hat. Also: daß unmöglich ist, Werke vom Glauben zu
scheiden, ja so unmöglich, als Brennen und Leuchten vom Feuer mag geschieden werden.“ — Wird dies Bollwerk
später vielleicht etwas ernster und scharfsichtiger angefochten, dann soll es an uns nicht fehlen, unser Panier:
nicht Arndt, sondern: nur Luther! auch in dieser Beziehung ernstlicher zu vertheidigen. Jetzt würden wir uns deß
schämen, unnöthigerweise hier Facta paradiren zu lassen.

Auf daß der Worte aber nicht zu viele werden, so will ich zum guten Schluß nur ein wenig von ferne auf unsere
übrigen Bollwerke wider den Pietismus — sei er nun Spenerisch oder Arndtisch — hinweisen. Sehet auf unsre
Kriege, die wir im Namen des HErrn führen. Wahr ist es bei uns, was Str. sagt, „unter harten Kämpfen hat sich
das Evangelium Bahn gemacht.“ Oder sollen wir ein wenig aufzählen? Unser Kämpfen und Streiten drehte sich um
Hauptartikel unseres theueren Bekenntnisses. Man betrachte nur ein wenig die 14 Jahrgänge des „Lutheraner“
und die 3 Jahrgänge von „Lehre und Wehre.“ Es galt den Artikel von der Rechtfertigung wider die Methodisten,
Jesuiten und Grabau; den von der Kirche und Amt
\quellseite{14}
wider Grabau und Consorten, sowie wider die Unirten; von den heiligen Sacramenten, von der Taufe wider die
Wiedertäufer, vom h. Abendmahl wider die Pseudolutheraner und alle Sacramentsschwärmer; den Artikel von der
h. Absolution wider die Methodisten, Reformirten \&c.; den Artikel von der Privatbeichte wider die unkirchlichen
Lutheraner; den Artikel vom Ansehen des Wortes wider grobe und feine Rationalisten; den Artikel von den letzten
Dingen wider die leidigen Chiliasten; den Artikel von der christlichen Freiheit wider die falschgeistlichen
Ceremonienhasser, die Kirchenregiments-Geiferer und den freiweltlichen Pöbel. Wer solches etliche Jahre mit
durchlebt und das Pulverriechen dadurch gewöhnt hat, daß ihm oft dabei nicht anders zu Muthe war, als sei der
„Lutheraner“ wider Alle und Alle wider ihn, der dichtet ihm das nicht an, er liebe pietistisch die Ruhe
oder stecke nach dem in Arndt’scher Richtung. Da höre nur Einer, welch’ ein Wetter entbrennt, wenn der
tapfere „Lutheraner“ (oder auch unsere sonstigen Organe) ein wenig die Flügel lüftet und die wider stolze,
hartnäckige Feinde der Wahrheit biblisch berechtigte, und im reformatorischen Zeitalter so mächtig erprobte,
heilige Waffe der Ironie ein klein wenig gebraucht. Da gibt’s Jammerlieder, ein Geheul und Gelage von hüben
und drüben, daß einem das Herz entfallen könnte — wenn’s nicht fest wäre. Ebräer 13, 8.

Sollte es noch Noth sein, nach Solchem auf noch sonderliche Specialia hinzuweisen? z.\,B. auf unsern Kampf wider
das Conventikelwesen, unsere Beschlüsse wider „das Vorbeten und Herzensgebet“ durch Laien im öffentlichen
Gottesdienst (siehe 4ten Synodalbericht von 1850); auf unser Bestreben, der Privatbeichte allhier zu ihrem
rechten kirchlichen Ansehen, sowie zu allgemeinerem Gebrauch zu verhelfen; auf unsere frei-evangelische Weise,
der altkirchlichen Gottesdienstordnung je mehr und mehr Eingang zu verschaffen; auf unser lediglich auf Gottes
Befehl sich gründendes, in dem Gebot der Nächstenliebe gefordertes Betreiben der Mission unter denen, in deren
Erbsitzen wir wohnen, dabei wir ganz und gar nicht in heut zu Tage beliebter Weise eine Ecclesiola in ecclesia
bauen, sondern als Kirche handeln, viel weniger dem allgemeinen marktschreierischen Herausspuffen der kleinsten
Kleinigkeiten huldigen; oder auf den Geist, der in unsern Katechismusschulen, in unsern Lehranstalten oder bei
den Versammlungen unserer Prediger und Lehrer auf Synoden und bei Conferenzen herrscht; oder auf das Unbehagen,
das deutsche Pietisten oder Arndtisten in allen unsern Gemeinden fühlen; auf unsere gut lutherische Saumseligkeit
hinsichtlich der Ausbreitung unserer Kirche u. s. w. — das stärkste Zeugniß wider jene erdichtete „erbauliche
Vielgeschäftigkeit“, welche aber größtentheils von dem entschiedenen Festhalten an der Lehre „vom Beruf“
herrührt und im Lande der unmäßigsten Speculation, des unsinnigsten Jagens und Haschens oft keine geringe
Anfechtung hat, um dem Strudel der Zeit zu widerstehen; oder auf das Zeugniß unserer Feinde, daß uns freispricht
von der Beschuldigung unserer Freunde?

Ob darum „der noch nicht ausgefochtene Amts- und Kirchenstreit“ wirklich
\quellseite{15}
bei uns „vor der Hand noch die wahre Lage der Dinge verdeckt“, das mag der geneigte Leser selbst
entscheiden. Die Zeit wird’s lehren. Freilich wird die Beantwortung dieser Frage: Luther oder Arndt? „keine
friedliche sein“. Weil dieser von Str. prophezeite Kampf und Streit wider die Arndtische Richtung im Grund doch
nur ein Streit wider den leidigen Pietismus ist, davon diese die Wurzel ist, so folgt es von selbst, daß der Streit
in seiner höchsten Blüthe die Ausfegung dieses Sauerteiges zum Ziel hat und „der Riß“ darum auch theilweise
„schmerzlicher und schädlicher, als in den bisherigen Kämpfen werden wird“; weil er aber nur die endliche
und letzte Aufklärung des trüben Handels ist, so wird nach unserer geringen Meinung — denn wir können uns
deß noch nicht rühmen, in und durch unser Lutherthum es bis zu der Gabe der Prophezie gebracht zu haben — doch
diesem Streit das Großartige und Thatenreiche des Hauptstreites wohl fehlen. Doch — die Erfahrung wird’s lehren!

\begin{unterschrift}
Ein schlichter einfältiger Prediger.
\end{unterschrift}
